% GENERATED FILE. Edit canonical lessons, then run npm run generate:books.
% canonical-source-hash: fnv1a64:89efd5d5764cb43c
% canonical-lessons: ES-C728-sabio, ES-C728-sabiduria, ES-C728-astuto, ES-C728-habil, ES-C728-ingenio

\chapter{Listos y sabios}
\label{ch:es-listos-y-sabios}
\hlchaptermodality{728}

\begin{chapteropening}
\textbf{By the end of this chapter:} \emph{I can describe someone wise, clever or skilful.}

Complete the last lesson of chapter 728: I can describe someone wise, clever or skilful.
\end{chapteropening}

\section[sabio]{sabio — wise; the wise person}

\label{lesson:ES-C728-sabio}



\begin{quote}
Before the new one: say the Spanish for polite, then the Spanish for rude.
\end{quote}

\subsection*{You'll want to know: sabio}
\textbf{sabio} — \textquotedblleft{}wise; the wise person\textquotedblright{}.

\begin{quote}
Mi abuelo era un hombre sabio.
\end{quote}

\begin{quote}
\emph{My grandfather was a wise man.}
\end{quote}

\emph{Un sabio} is a sage or a scholar.

\begin{cousinweb}
\emph{sabio} comes from Latin \emph{sapidus} (EtymDB), from \emph{sapere} \textquotedblleft{}to taste, to be wise\textquotedblright{}, the source of \emph{saber}.
\end{cousinweb}

\begin{grammarlens}[title={What you've built}]
One more word for telling, explaining and describing at B1.
\end{grammarlens}

\subsection*{Your turn}
\begin{itemize}
\raggedright
  \item \emph{Say it:} \emph{sabio}
  \item \emph{Say it:} \emph{sabio}, once more
  \item \emph{Say it:} a sentence of your own with \emph{sabio}
  \item \emph{Recall:} say the Spanish for polite, then the Spanish for rude, then say \emph{sabio} again
\end{itemize}

\begin{tcolorbox}[breakable,colback=teal!4,colframe=teal!35!black,title={Before you move on}]
What does sabio mean? (\textquotedblleft{}Wise; the wise person\textquotedblright{}.) Say it once more.
\end{tcolorbox}
\section[la sabiduría]{la sabiduría — wisdom}

\label{lesson:ES-C728-sabiduria}



\begin{quote}
Before the new one: say the Spanish for rude, then the Spanish for wise.
\end{quote}

\subsection*{You'll want to know: la sabiduría}
\textbf{la sabiduría} — \textquotedblleft{}wisdom\textquotedblright{}.

\begin{quote}
La sabiduría viene con los años.
\end{quote}

\begin{quote}
\emph{Wisdom comes with the years.}
\end{quote}

\begin{cousinweb}
\emph{sabiduría} belongs with \emph{sabio} \textquotedblleft{}wise\textquotedblright{} (EtymDB), and its form is usually traced to old \emph{sabidor} \textquotedblleft{}one who knows\textquotedblright{}, built on \emph{saber}, from Latin \emph{sapere} \textquotedblleft{}to taste, to be wise\textquotedblright{}.
\end{cousinweb}

\begin{grammarlens}[title={What you've built}]
One more word for telling, explaining and describing at B1.
\end{grammarlens}

\subsection*{Your turn}
\begin{itemize}
\raggedright
  \item \emph{Say it:} \emph{la sabiduría}
  \item \emph{Say it:} \emph{la sabiduría}, once more
  \item \emph{Say it:} a sentence of your own with \emph{la sabiduría}
  \item \emph{Recall:} say the Spanish for rude, then the Spanish for wise, then say \emph{la sabiduría} again
\end{itemize}

\begin{tcolorbox}[breakable,colback=teal!4,colframe=teal!35!black,title={Before you move on}]
What does la sabiduría mean? (\textquotedblleft{}Wisdom\textquotedblright{}.) Say it once more.
\end{tcolorbox}
\section[astuto]{astuto — clever, cunning}

\label{lesson:ES-C728-astuto}



\begin{quote}
Before the new one: say the Spanish for wise, then the Spanish for wisdom.
\end{quote}

\subsection*{You'll want to know: astuto}
\textbf{astuto} — \textquotedblleft{}clever, cunning\textquotedblright{}.

\begin{quote}
El zorro es astuto.
\end{quote}

\begin{quote}
\emph{The fox is cunning.}
\end{quote}

\begin{cousinweb}
\emph{astuto} is a learned borrowing from Latin \emph{astūtus} \textquotedblleft{}crafty\textquotedblright{} (EtymDB), built on \emph{astus} \textquotedblleft{}craft, cunning\textquotedblright{}.
\end{cousinweb}

\begin{grammarlens}[title={What you've built}]
One more word for telling, explaining and describing at B1.
\end{grammarlens}

\subsection*{Your turn}
\begin{itemize}
\raggedright
  \item \emph{Say it:} \emph{astuto}
  \item \emph{Say it:} \emph{astuto}, once more
  \item \emph{Say it:} a sentence of your own with \emph{astuto}
  \item \emph{Recall:} say the Spanish for wise, then the Spanish for wisdom, then say \emph{astuto} again
\end{itemize}

\begin{tcolorbox}[breakable,colback=teal!4,colframe=teal!35!black,title={Before you move on}]
What does astuto mean? (\textquotedblleft{}Clever, cunning\textquotedblright{}.) Say it once more.
\end{tcolorbox}
\section[hábil]{hábil — skilful, clever}

\label{lesson:ES-C728-habil}



\begin{quote}
Before the new one: say the Spanish for wisdom, then the Spanish for clever.
\end{quote}

\subsection*{You'll want to know: hábil}
\textbf{hábil} — \textquotedblleft{}skilful, clever\textquotedblright{}.

\begin{quote}
Es muy hábil con las manos.
\end{quote}

\begin{quote}
\emph{He is very good with his hands.}
\end{quote}

\emph{Día hábil} — a working day.

\begin{cousinweb}
\emph{hábil} is a learned borrowing from Latin \emph{habilis} \textquotedblleft{}apt, skilful\textquotedblright{} (EtymDB), from \emph{habēre} \textquotedblleft{}to hold\textquotedblright{}; English \emph{able} is its inherited French form.
\end{cousinweb}

\begin{grammarlens}[title={What you've built}]
One more word for telling, explaining and describing at B1.
\end{grammarlens}

\subsection*{Your turn}
\begin{itemize}
\raggedright
  \item \emph{Say it:} \emph{hábil}
  \item \emph{Say it:} \emph{hábil}, once more
  \item \emph{Say it:} a sentence of your own with \emph{hábil}
  \item \emph{Recall:} say the Spanish for wisdom, then the Spanish for clever, then say \emph{hábil} again
\end{itemize}

\begin{tcolorbox}[breakable,colback=teal!4,colframe=teal!35!black,title={Before you move on}]
What does hábil mean? (\textquotedblleft{}Skilful, clever\textquotedblright{}.) Say it once more.
\end{tcolorbox}
\section[el ingenio]{el ingenio — wit, ingenuity}

\label{lesson:ES-C728-ingenio}



\begin{quote}
Before the new one: say the Spanish for clever, then the Spanish for skilful.
\end{quote}

\subsection*{You'll want to know: el ingenio}
\textbf{el ingenio} — \textquotedblleft{}wit, ingenuity\textquotedblright{}.

\begin{quote}
Tiene mucho ingenio para resolver problemas.
\end{quote}

\begin{quote}
\emph{She is very ingenious at solving problems.}
\end{quote}

\emph{Un ingenio} can also be a clever device; in Latin America \emph{un ingenio} is a sugar mill.

\begin{cousinweb}
\emph{ingenio} was borrowed from Latin \emph{ingenium} \textquotedblleft{}inborn nature, talent\textquotedblright{} (EtymDB), built on \emph{in} and the root of \emph{gignere} \textquotedblleft{}to beget\textquotedblright{}: the gifts you are born with.
\end{cousinweb}

\begin{grammarlens}[title={What you've built}]
One more word for telling, explaining and describing at B1.
\end{grammarlens}

\subsection*{Your turn}
\begin{itemize}
\raggedright
  \item \emph{Say it:} \emph{el ingenio}
  \item \emph{Say it:} \emph{el ingenio}, once more
  \item \emph{Say it:} a sentence of your own with \emph{el ingenio}
  \item \emph{Recall:} say the Spanish for clever, then the Spanish for skilful, then say \emph{el ingenio} again
\end{itemize}

\begin{tcolorbox}[breakable,colback=teal!4,colframe=teal!35!black,title={Before you move on}]
What does el ingenio mean? (\textquotedblleft{}Wit, ingenuity\textquotedblright{}.) Say it once more.
\end{tcolorbox}
